% Auto-generated by experiments/19_report_visualization_enrichment.py
\section*{Visual Evidence Gallery}
\begin{figure}[!htbp]
\centering
\begin{subfigure}[t]{0.48\linewidth}
\centering
\includegraphics[width=\linewidth]{results/figures/report_enrichment/stage03_boundary_greeks_diagnostics_put_boundary_curve.pdf}
\caption{The put early-exercise boundary is extracted from low-spot exercise to continuation transitions in the CN/PSOR continuation premium.}
\end{subfigure}
\hfill
\begin{subfigure}[t]{0.48\linewidth}
\centering
\includegraphics[width=\linewidth]{results/figures/report_enrichment/stage03_boundary_greeks_diagnostics_call_boundary_curve.pdf}
\caption{The dividend-paying call boundary uses the opposite high-spot exercise orientation and is evaluated separately from put boundaries.}
\end{subfigure}
\end{figure}

\begin{figure}[!htbp]
\centering
\begin{subfigure}[t]{0.48\linewidth}
\centering
\includegraphics[width=\linewidth]{results/figures/report_enrichment/stage03_boundary_greeks_diagnostics_put_delta_surface.pdf}
\caption{Delta surfaces are interpreted together with masks around nonsmooth regions.}
\end{subfigure}
\hfill
\begin{subfigure}[t]{0.48\linewidth}
\centering
\includegraphics[width=\linewidth]{results/figures/report_enrichment/stage03_boundary_greeks_diagnostics_put_gamma_surface.pdf}
\caption{Gamma spikes near kinks and free boundaries are shown as diagnostic limitations rather than hidden failures.}
\end{subfigure}
\end{figure}

\FloatBarrier
